% ===========================================================================
%  板块五 · 含条件的最值（卡片 5 + 题 14–17）
% ===========================================================================
\section{含条件的最值}

\begin{strand}
\keyword{本板块主线}：用条件中的定值「1」沟通条件与目标。
\end{strand}

\block{知识精讲：三种处理方式}

\textbf{其一：乘「1」法（代换法）。}条件给出「和为 1」（或可化为和为定值），目标为倒数和时，把目标乘上这个「1」：
\[
a+b=1:\quad \frac1a+\frac1b=\left(\frac1a+\frac1b\right)(a+b)=2+\frac ba+\frac ab\ge 2+2=4.
\]
展开后常数项保留，互倒对用基本不等式处理。等号：$a=b=\tfrac12$，代回检验。（题 7 已用过三元版本。此法为何有效，见拓展 C：它在配平次数。）

\textbf{其二：和积互化（解不等式法）。}条件把 $a+b$ 与 $ab$ 联系在一起（如 $ab=a+b+3$）时，用 $a+b\ge 2\sqrt{ab}$ 或 $ab\le\left(\dfrac{a+b}2\right)^2$ 把其中一个换成另一个，得到单变量的不等式，解出范围。

\textbf{其三：双换元。}目标中两个分母各是 $x$、$y$ 的组合、无法各自化简时，设两个分母为新变量 $m$、$n$，反解出 $x$、$y$，条件与目标往往同时化简（题 17）。

\textbf{识别变形。}出题时常把「和为定值」隐藏起来：$\dfrac1a+\dfrac{2}{1-a}$ 即 $a+b=1$ 时的 $\dfrac1a+\dfrac2b$（令 $b=1-a$）；$\dfrac{1}{a+1}+\dfrac{1}{b+1}$ 中 $(a+1)+(b+1)=3$ 也是定值。先找出定值，再选方法。

\begin{problem}{题 14 （「1」的代换）}
已知 $a,b>0$ 且 $a+b=1$。

(1) 求 $\dfrac1a+\dfrac1b$ 的最小值；

(2) 求 $\dfrac1a+\dfrac2b$ 的最小值；

(3) 求 $\dfrac1a+\dfrac{2}{1-a}$（$0<a<1$）的最小值；

(4) \starred\ 求 $\dfrac{1}{a+1}+\dfrac{1}{b+1}$ 的最小值。
\hint{(2) 仍乘 $(a+b)$，展开后不互倒的两项同样可用基本不等式；(3) 先辨认它与哪一小问相同；(4) 令 $u=a+1$、$v=b+1$，$u+v$ 是定值吗？}
\end{problem}

\begin{problem}{题 15 \starred（和与积的互化）}
已知 $a,b>0$，且 $ab=a+b+3$，分别求 $ab$ 与 $a+b$ 的取值范围。
\hint{求 $ab$：用 $a+b\ge 2\sqrt{ab}$ 替换条件中的 $a+b$，设 $t=\sqrt{ab}$，移项后因式分解，注意其中一个因式恒正。求 $a+b$：改用「积 $\le$ 半和的平方」。端点能否取到，需回代检验（$a=b=3$）。}
\end{problem}

\begin{problem}{题 16 （乘「1」法）}
已知 $x,y>0$，且 $\dfrac1x+\dfrac9y=1$，求 $x+y$ 的最小值。
\hint{条件本身就是「1」，用 $x+y$ 乘它。与题 14(2) 对照：一处是「和为 1」，一处是「倒数和为 1」，方向相反。}
\end{problem}

\begin{problem}{题 17 \starred\starred（双换元）}
已知实数 $x,y$ 满足 $x>y>0$，且 $x+y\le 2$，求 $\dfrac{2}{x+3y}+\dfrac{1}{x-y}$ 的最小值。（书例 1.1-6，解答由讲义补全）
\hint{设 $m=x+3y$，$n=x-y$，先反解 $x$、$y$，看条件 $x+y\le 2$ 变成什么。目标 $\dfrac2m+\dfrac1n$ 没有「等于定值」的条件时，先乘 $(m+n)$ 再除回去，注意不等号的方向。}
\end{problem}

\clearpage
